\begin{center}\begin{tikzpicture}[font=\small,>=stealth,line join=round]
\draw[gray!40,rounded corners=3pt] (0,0) rectangle (14.2,8.3);\node at (7.1,7.700000000000001) {例4：负倒数型必须保留跳跃，不能把断点连起来};
\begin{scope}[shift={(1.205263157894737,3.1085365853658535)},x=0.5263157894736843cm,y=1.1219512195121952cm]\draw[->,gray] (-0.20000000000000018,0)--(11.2,0) node[right]{$x$};\draw[blue,thick] (1,2)--(4,2);\draw[blue,fill=white] (4,2) circle(2.7pt);\fill[blue] (1,2) circle(2.3pt);\node[below=7pt] at (1,-1.25){$1$};\draw[orange!85!black,thick] (4,-0.5)--(7,-0.5);\draw[orange!85!black,fill=white] (7,-0.5) circle(2.7pt);\fill[orange!85!black] (4,-0.5) circle(2.3pt);\node[below=7pt] at (4,-1.25){$4$};\draw[blue,thick] (7,2)--(10,2);\draw[blue,fill=white] (10,2) circle(2.7pt);\fill[blue] (7,2) circle(2.3pt);\node[below=7pt] at (7,-1.25){$7$};\fill[orange!85!black] (10,-0.5) circle(2.3pt);\node[below=7pt] at (10,-1.25){$10$};\end{scope}\node[anchor=west,align=left,text width=5cm] at (8.4,4){每段长度3\\[9pt]$f(1)=2,\ f(4)=-\frac12$\\[9pt]$f(7)=2$\\[9pt]$f(x+3)=-\frac1{f(x)}$\\[9pt]$2026=6\times337+4$};
\node[text width=12.6cm,align=center] at (7.1,.65) {每段左闭右开；按长度6向两侧重复，此处只画一部分。};
\end{tikzpicture}\end{center}
